\documentclass[../../../../SoftwareRequirementsSpecification.tex]{subfiles}
\begin{document}
% Richiede le macro definite in src/macros/useCaseMacros.tex
% (incluse nel preambolo di SoftwareRequirementsSpecification.tex)

\ucstart
\begin{xltabular}{\textwidth}{|l|X|}
  \hline
  \ucid{PT-08} \\ \hline
  \ucname{Visualizzazione Atleti} \\ \hline
  \ucactor{PT} \\ \hline
  \ucdescription{Il PT visualizza l'elenco dei propri
    atleti e ne consulta i dettagli.} \\ \hline
  \ucprecond{Il PT deve essere autenticato.} \\ \hline
  \ucpostcond{Il PT visualizza i dettagli dell'atleta
    selezionato.} \\ \hline
  \ucmainflow{
    \item Il PT apre la pagina dell'elenco dei propri atleti.
    \item Il sistema mostra l'elenco degli atleti seguiti dal PT, in
      ordine alfabetico per cognome. Per ciascun atleta mostra nome,
      cognome ed email.
    \item Il PT preme su un atleta dell'elenco.
    \item Il sistema mostra la pagina di dettaglio dell'atleta, con:
      \begin{itemize}
        \item i dati anagrafici (nome, cognome, età, email);
        \item i dati fisici (altezza e peso);
        \item le palestre in cui l'atleta si allena;
        \item l'elenco delle schede assegnate, con nome, periodo di
          validità e stato. Il sistema segnala le assegnazioni ancora
          in bozza (vedi Nota 5 del caso d'uso PT-05);
        \item l'elenco degli appuntamenti e delle richieste
          di appuntamento dell'atleta, con data, ora e stato.
      \end{itemize}
    \item Il PT può premere il pulsante "Assegna nuova Scheda" per
      assegnare una nuova scheda all'atleta.
    \item Il flusso continua dal passo 3 del caso d'uso PT-05.
  } \\ \hline
  \ucaltflow{
    \textbf{2a. Nessun atleta presente} \newline
    - Il sistema mostra un messaggio ("Nessun atleta registrato"). \newline
    - Il flusso termina. \newline
    \textbf{5a. Personalizzazione di una scheda già assegnata} \newline
    - Invece di assegnare una nuova scheda, il PT preme su una scheda
    confermata dell'elenco delle schede assegnate, per modificarne le
    sessioni di allenamento limitatamente a quell'atleta. \newline
    - Il flusso continua dal passo 3 del caso d'uso PT-09. \newline
    \textbf{5b. Ripresa di un'assegnazione in bozza} \newline
    - Il PT preme su un'assegnazione segnalata come bozza, per
    completarne i carichi. \newline
    - Il flusso continua dal flusso alternativo 2a del caso
    d'uso PT-05.
  } \\ \hline
  \ucnote{
    \textbf{Nota 1:} tutte le informazioni della pagina di dettaglio
    sono in sola lettura. Lo statement assegna al PT solo la
    visualizzazione dei dettagli dei propri atleti. L'inserimento e
    la modifica dei dati anagrafici spettano invece all'Atleta (caso
    d'uso A-04).\newline
    \textbf{Nota 2:} gli elenchi di schede, appuntamenti e richieste
    possono essere vuoti. In tal caso il sistema mostra la sezione
    vuota, senza cambiare il flusso.\newline
    \textbf{Nota 3:} il dettaglio mostra l'elenco delle schede
    assegnate anche per un altro motivo: permette al PT di
    verificare, prima di assegnarne una nuova con PT-05, che gli
    intervalli temporali delle schede dell'atleta non si
    sovrappongano.\newline
    \textbf{Nota 4:} il PT vede solo i propri atleti. L'elenco
    corrisponde agli atleti associati al PT autenticato.
  } \\ \hline
\end{xltabular}
\end{document}
